\subsubsection{Quantum Oracle Construction}\label{sec:quantum-oracle-construction}

This section should answer one question: ``\emph{\textbf{how do we build a quantum oracle that computes $f_w (x) = w \cdot x \mod 2$?}}''. The great thing is that the construction is extremely simple. For \textbf{each 1 in the hidden string, there is one CNOT gate} in the quantum circuit. Let's see how this works.

\highspace
The Bernstein-Vazirani algorithm assumes access to a quantum oracle implementing the function:
\begin{equation*}
    f_w (x) = w \cdot x \mod 2
\end{equation*}
As in Deutsch's algorithm, the oracle must be represented by a \textbf{unitary operator}, since every quantum computation must be reversible.

\highspace
The oracle acts on \textbf{two quantum registers}:
\begin{itemize}
    \item An $n$-\textbf{qubit data register} containing the input \ket{x}.
    \item A single \textbf{ancilla qubit} \ket{y}.
\end{itemize}
Its action is defined as follows:
\begin{equation}
    U_{f_w} \ket{x} \ket{y} = \ket{x} \ket{y \oplus f_w (x)}
\end{equation}
Notice that the \textbf{oracle never changes the input register}. It simply computes the function $f_w(x)$ and \textbf{XORs the result into the ancilla} qubit \ket{y}. This is exactly the same reversible construction used in Deutsch's algorithm, but now the input consists of $n$ qubits instead of one.

\highspace
\begin{flushleft}
    \textcolor{Green3}{\faIcon{lock} \textbf{Encoding the Hidden Bitstring}}
\end{flushleft}
The hidden string (algorithm's goal):
\begin{equation*}
    w = \left(w_1, w_2, \ldots, w_n\right)
\end{equation*}
Determines the oracle completely. Recall that the function $f_w(x)$ is defined as:
\begin{equation*}
    f_w (x) = w \cdot x \mod 2 \implies w_1 x_1 \oplus w_2 x_2 \oplus \ldots \oplus w_n x_n
\end{equation*}
\hl{Each term contributes to the final XOR only when the corresponding bit of $w$ is equal to $1$} (if $w_i = 0$, then $w_i x_i = 0$ regardless of the value of $x_i$). Therefore, the \textbf{oracle's action is determined entirely by the positions of the 1s in the hidden string}.

\highspace
This observation leads to a remarkably \textbf{simple implementation}. For every position $i$:
\begin{itemize}
    \item If $w_i = 1$, then insert a CNOT whose:
    \begin{itemize}
        \item \textbf{Control} qubit is the \textbf{data qubit} \ket{q_i} corresponding to the $i$-th bit of the input \ket{x}.
        \item \textbf{Target} qubit is the \textbf{ancilla qubit} \ket{y}.
    \end{itemize}
    \item If $w_i = 0$, then do nothing.
\end{itemize}
Therefore, \textbf{every 1 in the hidden string corresponds to exactly one CNOT gate}. The oracle is simply a collection of CNOT gates, one for each 1 in the hidden string.

\highspace
\begin{flushleft}
    \textcolor{Green3}{\faIcon{question-circle} \textbf{Why Does This Work?}}
\end{flushleft}
Recall that a CNOT flips its target only when its control qubit is equal to 1. Suppose the ancilla initially contains \ket{y}. If a control qubit has value $x_i = 1$, the corresponding CNOT performs the operation:
\begin{equation*}
    y \to y \oplus 1
\end{equation*}
If instead the control qubit has value $x_i = 0$, the ancilla is left unchanged:
\begin{equation*}
    y \to y \oplus 0 = y
\end{equation*}
Applying all the required CNOT gates one after another accumulates the XOR of all selected input bits into the ancilla qubit. The final value of the ancilla is:
\begin{equation*}
    y \to y \oplus \left(w_1 x_1\right) \oplus \left(w_2 x_2\right) \oplus \ldots \oplus \left(w_n x_n\right)
\end{equation*}
Which is exactly the desired result:
\begin{equation*}
    y \to y \oplus f_w (x)
\end{equation*}
Thus the oracle realizes the unitary transformation:
\begin{equation*}
    U_{f_w} \ket{x} \ket{y} = \ket{x} \ket{y \oplus f_w (x)}
\end{equation*}
For example, if the hidden string is $w = 10110$, the result is:
\begin{center}
    \begin{tabular}{@{} c c c @{}}
        \toprule
        \textbf{Position} & $w_i$ & \textbf{Oracle Action} \\
        \midrule
        1 & 1 & CNOT with control \ket{q_1} and target \ket{y} \\
        2 & 0 & Do nothing \\
        3 & 1 & CNOT with control \ket{q_3} and target \ket{y} \\
        4 & 1 & CNOT with control \ket{q_4} and target \ket{y} \\
        5 & 0 & Do nothing \\
        \bottomrule
    \end{tabular}
\end{center}
\noindent
Without knowing $w$, we do not know which CNOTs are present inside the oracle. The \textbf{purpose of the Bernstein-Vazirani algorithm is precisely to discover this hidden pattern} using a single oracle query. When we will introduce the connection with phase kickback, it will become clear how this is possible.

\highspace
\begin{flushleft}
    \textcolor{Green3}{\faIcon{link} \textbf{Connection with Phase Kickback}}
\end{flushleft}
At this point, the oracle may appear to do nothing more than compute a classical Boolean function. However, the crucial idea is that the ancilla will \textbf{not} be initialized in \ket{0}. Instead, it will be prepared in the state \ket{-}:
\begin{equation*}
    \ket{-} = \frac{\ket{0} - \ket{1}}{\sqrt{2}}
\end{equation*}
Which is an eigenstate of the Pauli-X operator with eigenvalue $-1$.

\highspace
Because each CNOT applies an $X$ gate to the ancilla whenever its control qubit is $1$, the oracle no longer stores the value of $f_w(x)$ in the ancilla. Instead, it transfers that information into a \textbf{phase} of the data register through \textbf{phase kickback}.

\highspace
This mechanism is the key to recovering the entire hidden bitstring with only one oracle query and will be examined in detail on \autopageref{sec:phase-kickback-intuition}.

\highspace
\begin{takeawaysbox}[: Quantum Oracle Construction]
    The Bernstein-Vazirani oracle is constructed by placing \textbf{one CNOT for every position where the hidden bitstring has a 1}, using the corresponding data qubit as the control and the ancilla as the target. Acting on \ket{x}\ket{y}, the oracle computes:
    \begin{equation*}
        U_{f_w} \ket{x} \ket{y} = \ket{x} \ket{y \oplus f_w (x)}
    \end{equation*}
    Where $f_w (x) = w \cdot x \mod 2$. When the ancilla is later prepared in \ket{-}, these CNOTs produce phase kickback instead of storing the function value directly.
\end{takeawaysbox}